\chapter{Week 7: Recursion and Objects}
\label{ch:week7}

Lecture 7: October 20, 2026. Recitation 7: October 23, 2026.

Assigned reading and resources (course page):
\begin{itemize}
  \item \cite{PLAI} Chapter 9 (Recursion and Cycles: Procedures and Data), which builds recursive functions out of a placeholder and mutation, as this chapter does.
  \item \cite{PLAI} Chapter 10 (Objects), whose objects-by-desugaring sections follow the same path as \Cref{sec:objects}: messages, constructors, state, static members, \code{self}, dynamic dispatch, and inheritance through a parent object.
  \item Lecture code: \code{lecture07.zip}. According to the course page, its interpreter that encodes recursion with boxes (\Cref{sec:recboxes}) is a bit cleaner than the version written in lecture.
\end{itemize}

Further reading:
\begin{itemize}
  \item \cite{TAPL} Section 5.2 (Programming in the Lambda-Calculus), whose last part on recursion defines \code{fix}, the call-by-value fixed-point combinator, which is \code{Z} of \Cref{sec:rec7}; optional.
\end{itemize}

Under lexical scope, a function cannot call itself by the name it is being defined with: when the closure is created, that name is not bound yet. This chapter adds a form \code{rec} for recursive functions and implements it twice. The first version teaches application to rebind the function's own name at every call. The second version uses no new kind of value at all. It builds the cycle that recursion needs out of a box, by first creating a placeholder and then filling it in, a technique called \emph{backpatching}, packaged as a function \code{fix}. The same ideas then build objects: an object is a closure that dispatches on messages, its fields are mutable variables or boxes, it finds \code{self} either by backpatching or by taking itself as an argument, and it inherits by delegating to a parent object. The recitation gets recursion with no mutation at all, from the Y and Z combinators. This chapter covers:
\begin{itemize}
  \item \Cref{sec:recproblem}: \code{if0} and \code{rec}, why \code{let} cannot define a recursive function, and a first implementation with a new value \code{recV};
  \item \Cref{sec:cycles}: cyclic data built from Racket boxes, backpatching a recursive function in Racket, and \code{fix};
  \item \Cref{sec:recboxes}: \code{rec} in the store-passing interpreter, desugared into a call to \code{fix}, or inlined with a fresh name, and why both need a forwarding function;
  \item \Cref{sec:objects}: objects as closures over messages, constructors, private, shared, and static state, two ways to get \code{self}, trees built from objects, and inheritance by delegation;
  \item \Cref{sec:rec7}: Recitation~7, circular lists and trees with boxes, fixed points, the Y and Z combinators, and \code{rec} implemented with \code{Z}.
\end{itemize}

\section{The Problem of Recursion}
\label{sec:recproblem}

\subsection{Two New Forms}

We start from the closure interpreter of \Cref{sec:closures} and add two forms. \code{(if0 e1 e2 e3)} evaluates \code{e1}. If the result is 0, it evaluates \code{e2}; otherwise it evaluates \code{e3}. This is enough branching to write a base case. \code{(rec f x e)} is a function with parameter \code{x} and body \code{e}, whose body may call the function itself under the name \code{f}.
\begin{lstlisting}
(define-type Expr
  [numC (n : number)]
  [plusC (e1 : Expr) (e2 : Expr)]
  [timesC (e1 : Expr) (e2 : Expr)]
  [letC (x : symbol) (e1 : Expr) (e2 : Expr)]
  [lambdaC (x : symbol) (e : Expr)]
  [appC (e1 : Expr) (e2 : Expr)]
  [idC (x : symbol)]
  [if0C (guard : Expr) (e1 : Expr) (e2 : Expr)]
  [recC (f : symbol) (x : symbol) (e : Expr)])
\end{lstlisting}
The parser gets one line for each:
\begin{lstlisting}
            [(rec) (recC (s-exp->symbol (second l)) (s-exp->symbol (third l)) (parse (fourth l)))]
            [(if0) (if0C (parse (second l)) (parse (third l)) (parse (fourth l)))]
\end{lstlisting}
\code{if0} only evaluates the branch it chooses, just like \code{if} in Racket:
\begin{lstlisting}
    [if0C (guard e1 e2)
          (let [(v (eval-env env guard))]
            (if (equal? v (numV 0))
                (eval-env env e1)
                (eval-env env e2)))]
\end{lstlisting}
\code{(if0 0 1 2)} is 1 and \code{(if0 1 1 2)} is 2. The guard's value is compared with \code{(numV 0)} as a whole, so a guard that is not a number, such as a closure, is simply ``not 0'' and selects \code{e3} rather than raising an error. With both forms, factorial is written
\begin{lstlisting}
(rec fact n (if0 n 1 (* n (fact (+ n -1)))))
\end{lstlisting}

\subsection{Why \texttt{let} Is Not Enough}

Without \code{rec}, the natural attempt is to bind the function with \code{let}:
\begin{lstlisting}
(let fact (lambda n (if0 n 1 (* n (fact (+ n -1)))))
  (fact 5))
\end{lstlisting}
Under lexical scope this fails with \code{No binding found}. The \code{lambda} is evaluated \emph{before} \code{let} binds \code{fact}, so the closure saves an environment without \code{fact}. When the body later calls \code{fact}, it looks the name up in that saved environment and does not find it.

Under dynamic scope, the same program works. The body is evaluated in the environment of the \emph{call} \code{(fact 5)}, and by then \code{let} has bound \code{fact}. This is one of the few things dynamic scope makes easier, and it is no reason to go back to it (\Cref{sec:dynscope}). We want recursion and lexical scope together.

\subsection{A First Solution: \code{recV}}

The function needs an environment that contains the function itself. That is a chicken-and-egg problem. To build the closure we need its environment first, and the environment is supposed to contain the closure.

The first solution avoids the cycle by \emph{waiting}. A new kind of value records the function's own name next to the usual closure fields:
\begin{lstlisting}
(define-type Value
  [numV (n : number)]
  [closV (env : Env) (x : symbol) (e : Expr)]
  [recV (env : Env) (f : symbol) (x : symbol) (e : Expr)])
\end{lstlisting}
\code{rec} only packages these fields, without binding \code{f}:
\begin{lstlisting}
    [recC (f x e) (recV env f x e)]
\end{lstlisting}
The binding happens at each call. Application now has to tell the kinds of function apart:
\begin{lstlisting}
    [appC (e1 e2)
          (let [(v1 (eval-env env e1))
                (v2 (eval-env env e2))]
            (type-case Value v1
              [closV (env x e)
                     (eval-env (extend-env (bind x v2) env) e)]
              [recV (env f x e)
                    (eval-env (extend-env (bind f v1) (extend-env (bind x v2) env)) e)]
              [numV (n)
                    (error 'eval-env "Error: applied number as a function")]))]
\end{lstlisting}
In the \code{recV} case, \code{v1} \emph{is} the function being called. So binding \code{f} to \code{v1} makes the function's own name visible in its body. By the time of the call, \code{v1} already exists, so there is no cycle to build. Every recursive call goes through \code{appC} again and adds the binding again, so recursion can go as deep as needed.

For \code{((rec fact n ...) 2)}, the body runs in
\begin{lstlisting}[style=repl]
[fact = recV(...), n = 2, the environment where rec was evaluated]
\end{lstlisting}
and the call \code{(fact (+ n -1))} in the body finds the same \code{recV}.

The order of the two bindings matters only when \code{f} and \code{x} are the same name. Here \code{f} is added last, so in \code{(rec g g ...)} the name \code{g} in the body means the function, not the argument. Binding \code{x} last instead would let the parameter shadow the function name, which is what Racket does for \code{(letrec ([g (lambda (g) ...)]) ...)}.

The lecture's examples:
\begin{lstlisting}
((rec fact x (if0 x 1 (* x (fact (+ x -1))))) 5)           ; 120

(let fact (rec f x (if0 x 1 (* x (f (+ x -1)))))
  (fact 6))                                                ; 720

(let fibb (rec f x (if0 x
                        0
                        (if0 (+ x -1) 1
                             (+ (f (+ x -1)) (f (+ x -2))))))
  (fibb 6))                                                ; 8
\end{lstlisting}
In the last two, the name the body uses for itself, \code{f}, is not the name the rest of the program uses, \code{fact} or \code{fibb}. The \code{let} only names the finished \code{recV} for the outside. The recursive calls go through \code{f}, which \code{appC} binds at every call.

\section{Cycles from Boxes}
\label{sec:cycles}

The \code{recV} solution needs a new kind of value and a special case in application. Racket's own \code{letrec} works differently: it builds the cycle directly, with mutation. This section shows the idea in plain \code{\#lang racket}.

\subsection{A Box That Contains Itself}

The direct attempt at a value that refers to itself is
\begin{lstlisting}
(let ([b b])
  b)
\end{lstlisting}
and it fails with \code{b: unbound identifier}, for the reason of \Cref{sec:recproblem}: the right-hand side is evaluated before \code{b} is bound. A box gets around this:
\begin{lstlisting}
(let ([b (box 'dummy)])
  (begin
    (set-box! b b)
    b))
\end{lstlisting}
The program first creates a box holding a placeholder, \code{'dummy}. It then replaces the contents with the box itself. The result is a cycle of length one: unboxing \code{b} gives \code{b} again, however many times we unbox. Racket prints it as
\begin{lstlisting}[style=repl]
#0='#&#0#
\end{lstlisting}
Here \code{\#\&} marks a box, \code{\#0=} labels a value as number 0, and \code{\#0\#} refers back to that value. So the output reads ``a box whose contents are value 0'', which is the box itself.

A cycle of length two works the same way:
\begin{lstlisting}
(let* ([b1 (box 'dummy)]
       [b2 (box b1)])
  (set-box! b1 b2)
  b1)
\end{lstlisting}
\code{b2} is built holding \code{b1}, which already exists. Then \code{b1} is changed to hold \code{b2}, which closes the cycle. Racket prints \code{\#0='\#\&\#\&\#0\#}, a box holding a box holding the first box. \code{let*} is needed because the initial value of \code{b2} refers to \code{b1}.

\begin{remark}
These examples only run in \code{\#lang racket}. In \code{plai-typed}, the body of a \code{let*} must be a single expression, so the second example is a syntax error without a \code{begin}. Worse, neither example can be typed: a box that contains itself would need the infinite type ``box of box of box of \dots''. The checker reports \code{(boxof symbol) vs.\ symbol} for the first one.
\end{remark}

\subsection{Backpatching}

Both examples follow the same three steps, and so does every cycle built this way:
\begin{enumerate}
  \item \emph{Placeholder}: create a box holding a dummy value.
  \item \emph{Build}: create the things that refer to the box. This is possible because the box already exists.
  \item \emph{Patch}: use \code{set-box!} to put the real value into the box, which closes the cycle.
\end{enumerate}
This technique is called \emph{backpatching}. It is exactly what a recursive function needs. The function's environment must refer to the function. With backpatching, the environment refers to a box instead, and the box is filled with the function once the function exists.

\subsection{A Recursive Function in Racket}

\begin{lstlisting}
(let* ([fact (box 'dummy)]
       [fact-fun (lambda (n)
                   (if (equal? n 0)
                       1
                       (* n ((unbox fact) (+ n -1)))))])
  (begin
    (set-box! fact fact-fun)
    ((unbox fact) 5)))
\end{lstlisting}
The three steps are visible. \code{fact} is the placeholder box. \code{fact-fun} is a closure whose saved environment contains \code{fact}. \code{set-box!} puts the closure into the box. From then on, the body finds the box through \code{fact}, unboxes it, and gets \code{fact-fun} itself. The result is 120.

\begin{pitfall}[\code{let} instead of \code{let*}]
With \code{let}, the same program is rejected with \code{fact: unbound identifier}. In a \code{let}, every right-hand side is evaluated in the outer scope, so the \code{lambda} cannot see the \code{fact} bound next to it (\Cref{sec:substrev}). This is the problem of \Cref{sec:recproblem} once more: the function's own name must be in scope when its closure is created.
\end{pitfall}

Note that the body has to write \code{((unbox fact) (+ n -1))}, not \code{(fact (+ n -1))}. The programmer has to know that \code{fact} is a box.

\subsection{\texttt{fix}: Backpatching Once and for All}

The three steps do not depend on what the function computes, so they can be written once, as a function \code{fix}. Its argument \code{f} is the function to make recursive, written so that it receives ``the recursive call'' as an extra first argument \code{rec}:
\begin{lstlisting}
(define fix
  (lambda (f)
    (let* ([b (box 'dummy)]
           [frec (lambda (x)
                   (f (unbox b) x))])
      (begin
        (set-box! b frec)
        frec))))

(define fact-alt
  (fix (lambda (rec n)
         (if (equal? 0 n)
             1
             (* n (rec (- n 1)))))))
\end{lstlisting}
\code{(fact-alt 5)} is 120. The steps are as before, with one change: the box holds \code{frec}, not \code{f}. \code{frec} takes the real argument \code{x}, fetches itself from the box, and calls \code{f} with itself as \code{rec}. So when \code{f}'s body calls \code{(rec (- n 1))}, it calls \code{frec} again, which goes through the box again. The function passed to \code{fix} never mentions a box: it calls \code{rec} like any other function. \Cref{sec:recboxes} does the same inside our interpreter, with \code{fix} written in the interpreted language.

\section{\texttt{rec} with Boxes}
\label{sec:recboxes}

\subsection{The Starting Point}

To use backpatching inside our own language, the language needs boxes. So the lecture starts again from the box interpreter of \Cref{sec:storepass} and adds \code{if0C} and \code{recC} to it. \tyname{Expr} has \code{boxC}, \code{unboxC}, \code{setboxC}, and \code{seqC}; \tyname{Value} has \code{boxV} but no \code{recV}; and \code{eval-env} takes a store and returns a \tyname{Result}:
\begin{lstlisting}
(define (eval-env (env : Env) (sto : Store) (e : Expr)) : Result
  ...)
\end{lstlisting}
Every case threads the store: the store returned by one subexpression is passed to the next one. As in \Cref{sec:storepass}, \code{let} is not a case of its own but is rewritten into an application:
\begin{lstlisting}
    [letC (x e1 e2) (eval-env env sto (appC (lambdaC x e2) e1))]
\end{lstlisting}

\subsubsection{Threading the store through \code{if0}} \code{if0} changes in the same way as \code{plusC}:
\begin{lstlisting}
    [if0C (guard e1 e2)
          (type-case Result (eval-env env sto guard)
            (res (v sto-1)
                 (if (equal? v (numV 0))
                     (eval-env env sto-1 e1)
                     (eval-env env sto-1 e2))))]
\end{lstlisting}
There are three reasons for this shape. First, \code{eval-env} now returns a \tyname{Result}, so the value \code{v} has to be taken out before it can be compared with 0. Second, the guard may have side effects, and the chosen branch must see them. So the branch runs in \code{sto-1}, not \code{sto}:
\begin{lstlisting}
(let b (box 5)
  (if0 (begin (set-box! b 0) (unbox b))
       (unbox b)    ; must see 0, the value the guard stored
       99))
\end{lstlisting}
With \code{sto}, the branch would still see 5. Third, the branch's own \tyname{Result} is returned as it is, so the store after the branch reaches the rest of the program.

\subsection{Desugaring \texttt{rec} into \texttt{fix}}

The Racket \code{fix} of \Cref{sec:cycles} can be written in our language too. Our functions take exactly one argument, so the function passed to \code{fix} takes the recursive call first and returns a function of \code{x}: it is \code{(lambda f (lambda x e))} instead of \code{(lambda (rec n) ...)}. Then
\begin{lstlisting}
(rec f x e)  =  (fix (lambda f (lambda x e)))
\end{lstlisting}
and \code{fix} is an ordinary program in our language, parsed once:
\begin{lstlisting}
(define fix
  (parse
   '(lambda f
      (let b (box 0)
        (let frec (lambda x ((f (unbox b)) x))
          (begin
            (set-box! b frec)
            frec))))))
\end{lstlisting}
The \code{recC} case builds the application and evaluates it:
\begin{lstlisting}
    [recC (f x e)
          (eval-env env sto (appC fix (lambdaC f (lambdaC x e))))]
\end{lstlisting}
\code{fix} is an \tyname{Expr}, not a Racket function: the interpreter splices it into the program as if the programmer had written it. The pieces around it are \tyname{Expr} constructors too, \code{(lambdaC f (lambdaC x e))} and not a \code{closV}, because \code{eval-env} evaluates them into closures itself. With this version, \code{recC} produces an ordinary \code{closV}, so the \code{recV} value of \Cref{sec:recproblem} is no longer needed. Both examples of \Cref{sec:recproblem} still give 720 and 8.

\subsubsection{What \code{fix} does} The steps are those of the Racket \code{fix}. \code{b} is the placeholder box, holding 0 because the language has no symbols. \code{frec} is created inside the scope of \code{b}. \code{set-box!} puts \code{frec} into \code{b}, and \code{fix} returns \code{frec}. When the program calls \code{frec} with an argument \code{v}:
\begin{enumerate}
  \item \code{(unbox b)} fetches \code{frec} itself from the store;
  \item \code{(f frec)} returns the real function \code{(lambda x e)}, closed over an environment in which \code{f} is \code{frec};
  \item that function is applied to \code{v}, and every call to \code{f} in \code{e} is a call to \code{frec}, which repeats from step 1.
\end{enumerate}
So the real function is rebuilt at every call; only \code{frec} exists once. The store after the factorial example shows the backpatch. It holds two cells for location 0, the newer one on top:
\begin{lstlisting}[style=repl]
(list (cell 0 (closV [b = (boxV 0), f = (closV '() 'f (lambdaC 'x ...))]
                     'x
                     ((f (unbox b)) x)))      ; frec, written by set-box!
      (cell 0 (numV 0)))                      ; the placeholder, written by box
\end{lstlisting}
\code{frec}'s environment contains \code{(boxV 0)}, and location 0 contains \code{frec}: this is the cycle, running through the store.

\subsubsection{No fresh names needed} The names \code{b} and \code{frec} are bound inside \code{fix}'s body, and \code{e} is not written inside that body. \code{e} is part of the argument \code{(lambda f (lambda x e))}, which is evaluated in the environment of the \code{rec}, before \code{fix} binds anything. By lexical scope, \code{fix}'s names are invisible to \code{e}, even when \code{e} uses the same names from outside:
\begin{lstlisting}[style=repl]
(let b 100 ((rec g n (if0 n b (g (+ n -1)))) 3))
---> 100
\end{lstlisting}
The \code{b} in the body is the outer \code{b}, not \code{fix}'s box.

\subsection{Inlining \texttt{fix}}

Instead of calling \code{fix}, \code{recC} can write the backpatching out in place. Writing \code{z} for a fresh name,
\begin{lstlisting}
(rec f x e)
  =  (let z (box 0)
       (let f (lambda x ((unbox z) x))
         (begin (set-box! z (lambda x e))
                f)))
\end{lstlisting}
The roles are swapped compared with \code{fix}. Here \code{f} is the forwarding function, and the box holds the real function \code{(lambda x e)}, which is created inside \code{(let f ...)} so its closure can see \code{f}. As a \code{recC} case:
\begin{lstlisting}
    [recC (f x e)
          (let ([z (gensym)])
            (eval-env env sto
                      (letC z (boxC (numC 0))
                            (letC f (lambdaC x (appC (unboxC (idC z)) (idC x)))
                                  (seqC (setboxC (idC z) (lambdaC x e))
                                        (idC f))))))]
\end{lstlisting}
Now \code{e} \emph{is} written inside the scope of \code{z}, so the box's name must be fresh. It comes from \code{gensym}, imported with \code{typed-in} as in \Cref{sec:capavoid}. A fixed name such as \code{'z} would shadow any \code{z} that \code{e} uses from outside:
\begin{lstlisting}[style=repl]
(let z 100 ((rec g n (if0 n z (g (+ n -1)))) 3))
---> 100        with gensym
---> a box      with the fixed name 'z
\end{lstlisting}

\begin{keypoint}[A closed helper needs no fresh names]
A desugaring that puts the user's code inside new binders has to invent fresh names for them. Passing the user's code as an \emph{argument} to a closed function such as \code{fix} avoids the problem: lexical scope keeps the helper's names and the user's names apart.
\end{keypoint}

\subsection{Why a Forwarding Function?}

In \code{fix}, \code{frec} is not the box, and not the result of \code{(f (unbox b))} either. It is a small function that does the unboxing when it is called. The two simpler choices both fail.

\paragraph{Return the box.} If \code{fix} returned \code{b} itself, the program would get a \code{boxV}, and a call such as \code{(fact 6)} would apply a box, which is an error. The body would have to write \code{((unbox f) (+ x -1))}, as the Racket version of \Cref{sec:cycles} does. But the point of \code{rec} is that the programmer calls \code{f} like any other function.

\paragraph{Unbox right away.} Write \code{(let frec (f (unbox b)) ...)} instead. Our \code{let} is eager, so \code{(unbox b)} runs at once, before \code{set-box!}. \code{f} receives the placeholder 0 as its recursive call, and the first recursive call in \code{e} applies the number 0. Patching the box later changes nothing, since 0 has already been passed.

\paragraph{The forwarding function.} \code{(lambda x ((f (unbox b)) x))} is a function, so the program can call it directly. Creating it does not run \code{(unbox b)}. The unbox runs only when \code{frec} is called, and by then the box has been patched. The inlined version's \code{(lambda x ((unbox z) x))} works for the same reason.

\begin{keypoint}[Delaying with a \code{lambda}]
Replacing an expression \code{g} by \code{(lambda x (g x))} is called \emph{$\eta$-expansion}. Both behave the same when called, but the second one evaluates \code{g} only at the call. This is the tool for any value that is not ready now but will be ready by the time it is used. The Z combinator of \Cref{sec:rec7} uses it to get recursion with no boxes at all.
\end{keypoint}

\begin{remark}
Two other designs avoid the forwarding function. One substitutes \code{(unbox z)} for \code{f} in \code{e}, which needs the capture-avoiding substitution of \Cref{sec:capavoid}. The other uses mutable variables, where the environment maps names to locations (\Cref{sec:rec6}): bind \code{f} to a placeholder, create the closure, and then \code{set!} \code{f} to it. That is how Racket's \code{letrec} behaves.
\end{remark}

\section{Objects from Closures}
\label{sec:objects}

This section asks whether objects need to be built into a language, or whether the features we already have are enough. The answer is that closures and mutable state go a long way. The code in this section is plain \code{\#lang racket}.

\subsection{An Object Is a Function on Messages}

An object answers \emph{messages}. In a language like Java, \code{o.add(x, y)} sends the message \code{add} with arguments \code{x} and \code{y} to \code{o}. We can model \code{o} as a function that takes a message name and returns the method for it:
\begin{lstlisting}
(define o-1
  (lambda (m)
    (case m
      [(add) (lambda (x y) (+ x y))]
      [(add1) (lambda (x) (+ x 1))]
      [(sub1) (lambda (x) (- x 1))])))
\end{lstlisting}
\code{((o-1 'add1) 5)} is 6 and \code{((o-1 'add) 5 6)} is 11. The \code{case} on the message is the object's \emph{dispatch}: it picks the method that the message names.

Writing two calls each time is clumsy, so a helper sends a message with any number of arguments:
\begin{lstlisting}
(define (msg o m . a)
  (apply (o m) a))
\end{lstlisting}
The dot in the parameter list makes \code{a} a \emph{rest parameter}: it collects all remaining arguments into a list. \code{apply} then calls the method with the elements of that list as its arguments. So \code{(msg o-1 'add 5 6)} is \code{((o-1 'add) 5 6)}, which is 11.

\begin{pitfall}[\code{(lambda m ...)} versus \code{(lambda (m) ...)}]
Without parentheses, \code{(lambda m body)} also takes any number of arguments and binds \code{m} to the \emph{list} of them. In an object written that way, \code{(o 'first)} binds \code{m} to \code{'(first)}, no \code{case} clause matches, and the call quietly returns nothing.
\end{pitfall}

\subsection{Constructors and Fields}

A \emph{constructor} is a function that returns an object. Its parameters become the object's fields, because the methods are closures that capture them:
\begin{lstlisting}
(define (o-constr-1 x)
  (lambda (m)
    (case m
      ((addX) (lambda (y) (+ x y))))))
\end{lstlisting}
\code{(msg (o-constr-1 10) 'addX 5)} is 15. Each call of \code{o-constr-1} creates a new closure with its own \code{x}, so each object has its own field.

\subsection{State}

Fields captured this way cannot change. For mutable fields, the constructor puts the value in a box, as in \Cref{sec:scopestate}:
\begin{lstlisting}
(define (obj-counter count)
  (let [(b (box count))]
    (lambda (m)
      (case m
        [(inc) (lambda () (set-box! b (+ (unbox b) 1)))]
        [(dec) (lambda () (set-box! b (- (unbox b) 1)))]
        [(get) (lambda () (unbox b))]))))
\end{lstlisting}
Every call of the constructor creates a new box, so two counters do not affect each other:
\begin{lstlisting}[style=repl]
(define cntr (obj-counter 0))
(define cntr2 (obj-counter 10))
(msg cntr 'inc)
(msg cntr 'inc)
(msg cntr 'dec)
(msg cntr 'get)     ; 1
(msg cntr2 'inc)
(msg cntr2 'get)    ; 11
(msg cntr 'get)     ; 1, unchanged by cntr2
\end{lstlisting}
The box is also \emph{private}. Only the methods can reach it, since nothing outside the closure has a name for \code{b}. The only way to change a counter is to send it \code{inc} or \code{dec}.

\subsubsection{Where the box is created decides who shares it} The next version moves the box out of the object:
\begin{lstlisting}
(define (obj-counter-share count)
  (let ([b (box count)])
    (lambda (count)
      (begin (set-box! b count)
             (lambda (m)
               (case m
                 [(inc) (lambda () (set-box! b (+ (unbox b) 1)))]
                 [(dec) (lambda () (set-box! b (- (unbox b) 1)))]
                 [(get) (lambda () (unbox b))]))))))
\end{lstlisting}
\code{obj-counter-share} now returns a \emph{constructor}, and all objects made by that constructor share one box:
\begin{lstlisting}[style=repl]
(define mk (obj-counter-share 0))
(define a (mk 1))
(define b (mk 10))      ; resets the shared box to 10
(msg a 'inc)
(list (msg a 'get) (msg b 'get))   ; '(11 11)
\end{lstlisting}
Two details matter. Objects share the box only if they come from the \emph{same} \code{mk}; a second call of \code{obj-counter-share} creates a second box. And every new object runs \code{(set-box! b count)}, which overwrites the value seen by all existing objects.

\subsubsection{Static fields} The usual way to share state among all objects of a class is to create the shared state once, when the constructor itself is defined. The lecture writes this version with Racket's \code{set!} instead of boxes:
\begin{lstlisting}
(define o-static-1
  (let ([numobj 0])
    (lambda (amount)
      (begin
        (set! numobj (+ 1 numobj))
        (lambda (m)
          (case m
            [(inc) (lambda (n) (set! amount (+ amount n)))]
            [(dec) (lambda (n) (set! amount (- amount n)))]
            [(get) (lambda () amount)]
            [(count) (lambda () numobj)]))))))
\end{lstlisting}
The \code{let} runs once, when \code{o-static-1} is defined, so \code{numobj} is shared by every object: it is a \emph{static field}, and here it counts the objects created. The parameter \code{amount} is bound anew on every constructor call, so it belongs to one object: it is an \emph{instance field}.
\begin{lstlisting}[style=repl]
(let ([o (o-static-1 1000)])
  (msg o 'count))         ; 1
(let ([o (o-static-1 0)])
  (msg o 'count))         ; 2
\end{lstlisting}
\code{set!} changes a variable, while \code{set-box!} changes a box that a variable names. The behavior is the same as with \code{(box 0)} for \code{numobj} and \code{(box amount)} for the counter: a Racket variable that is assigned is in effect a box that the language creates for us. In the terms of \Cref{sec:rec6}, the environment maps the name to a location, and \code{set!} overrides that location.

\begin{keypoint}[Lifetime of a binding = how widely it is shared]
State created outside the constructor is shared by all objects; state created inside it, including the constructor's own parameters, belongs to one object. Static versus instance fields is just the question of which \code{let} or \code{lambda} the variable is bound by.
\end{keypoint}

\subsection{\texttt{self}}

A method often needs to send a message to its own object, as \code{this.second(x)} does in Java. The object therefore needs a name for itself. This is the problem of \Cref{sec:recproblem} again: the object must refer to the object, but it does not exist until it is built. There are two solutions.

\subsubsection{\code{self} through a box} The first uses the backpatching of \Cref{sec:cycles}:
\begin{lstlisting}
(define o-self!
  (let ([self (box 'dummy)])
    (begin
      (set-box! self
            (lambda (m)
              (case m
                [(first) (lambda (x) (msg (unbox self) 'second (+ x 1)))]
                [(second) (lambda (x) (+ x 1))])))
      (unbox self))))
\end{lstlisting}
The three steps are the same as before. \code{self} is a placeholder box. The object is built inside its scope, so the method \code{first} can reach the box. \code{set-box!} puts the object into the box. When \code{first} runs, \code{(unbox self)} is the object itself. \code{(msg o-self! 'first 5)} computes \code{5 + 1}, sends \code{second}, and returns 7. The name \code{self} plays exactly the role that \code{f} played for \code{rec}.

\subsubsection{\code{self} as an argument} The second solution needs no mutation. Every method takes the object as an extra first argument:
\begin{lstlisting}
(define o-self
  (lambda (m)
    (case m
      [(first) (lambda (self x) (msg/self self 'second (+ x 1)))]
      [(second) (lambda (self x) (+ x 1))])))
\end{lstlisting}
The sender passes it, through a variant of \code{msg} that supplies the object twice, once to find the method and once as its first argument:
\begin{lstlisting}
(define (msg/self o m . a)
  (apply (o m) o a))
\end{lstlisting}
\code{(msg/self o-self 'first 5)} is 7, as before. This is the explicit \code{self} parameter of Python methods: \code{o.f(x)} means \code{f(o, x)}. Every call between methods must then go through \code{msg/self}; a plain \code{msg} would pass one argument too few.

\begin{center}
\begin{tabular}{@{}lll@{}}
\toprule
& \code{o-self!} & \code{o-self} \\
\midrule
how a method finds \code{self} & through a box in its closure & the caller passes it \\
needs mutation & yes (backpatching) & no \\
send with & \code{msg} & \code{msg/self} \\
\bottomrule
\end{tabular}
\end{center}

\subsubsection{Messages are values} Since a message is just a symbol, the program can compute which message to send:
\begin{lstlisting}
(define test
  (lambda (x)
    (msg o-self! (if (equal? x 1) 'first 'second) 5)))
\end{lstlisting}
\code{(test 1)} sends \code{first} and gives 7; \code{(test 2)} sends \code{second} and gives 6. The method that runs is chosen at run time from a value. An ordinary method call such as \code{o.first(5)} fixes the method name in the program text; with closures, the name is just data.

\subsection{Objects with Structure: Trees}

Objects can refer to other objects. A binary tree is a node object whose fields are a value and two subtrees, and an empty-tree object \code{mt} with no fields:
\begin{lstlisting}
(define (mt)
  (lambda (m)
    (case m
      [(add) (lambda () 0)])))

(define (node v l r)
  (lambda (m)
    (case m
      [(add) (lambda () (+ v
                           (msg l 'add)
                           (msg r 'add)))])))
\end{lstlisting}
\code{add} sums the values in the tree. A node answers it by sending \code{add} to its two subtrees and combining the answers; the empty tree answers 0. \code{mt} is a constructor with no arguments, so an empty tree is written \code{(mt)}.
\begin{lstlisting}
(define a-tree
  (node 10
        (node 5 (mt) (mt))
        (node 15 (node 6 (mt) (mt)) (mt))))
\end{lstlisting}
\code{(msg a-tree 'add)} is $10 + 5 + 15 + 6 = 36$.

The node never asks whether a subtree is a node or an empty tree. It sends the message, and each subtree runs its own method. Which code runs is decided by the object that receives the message, not by the sender. This is \emph{dynamic dispatch}.

\subsection{Inheritance by Delegation}

Now add a \code{size} message, which counts the nodes, without editing \code{node} and \code{mt}. A new constructor takes the old one as an argument, the \emph{parent}, and builds a parent object with the same fields. It answers \code{size} itself and passes every other message on to the parent object:
\begin{lstlisting}
(define (node/size parent-maker v l r)
  (let ([parent-object (parent-maker v l r)])
    (lambda (m)
          (case m
            [(size) (lambda () (+ 1
                                  (msg l 'size)
                                  (msg r 'size)))]
            [else (parent-object m)]))))

(define (mt/size parent-maker)
  (let ([parent-object (parent-maker)])
    (lambda (m)
      (case m
        [(size) (lambda () 0)]
        [else (parent-object m)]))))
\end{lstlisting}
The \code{else} clause is the whole mechanism: a message the child does not know is \emph{delegated} to the parent, so the child inherits \code{add}. Every node of the tree has to be built with the new constructors, with \code{node} and \code{mt} passed as the parents:
\begin{lstlisting}
(define a-tree/size
  (node/size node
             10
             (node/size node 5 (mt/size mt) (mt/size mt))
             (node/size node 15
                        (node/size node 6 (mt/size mt) (mt/size mt))
                        (mt/size mt))))
\end{lstlisting}
\code{(msg a-tree/size 'size)} is 4, and \code{(msg a-tree/size 'add)} is still 36. For \code{add}, the outer \code{node/size} object delegates to its parent \code{node} object. That object sends \code{add} to \code{l} and \code{r}, which are \code{node/size} objects, and they delegate in turn. All subtrees must be \code{/size} objects: if one were a plain \code{node}, \code{size} would reach it and match no \code{case} clause. The dispatch would return \code{\#<void>} instead of a method, and \code{msg} would fail with \code{application: not a procedure}.

\begin{remark}
The parent object is built from \code{v}, \code{l}, and \code{r} alone; it never learns that a child exists. So delegation passes on the message but not the object. If a parent method sent a message to itself through a \code{self} of its own, as \code{o-self!} does, the message would reach the parent, and a method that the child overrides would not be used. Inheritance in the usual sense makes \code{self} mean the child, for example by passing the child as \code{self} in the style of \code{o-self}.
\end{remark}

\section{Recitation 7: Recursion}
\label{sec:rec7}

This recitation asks two closely related questions: how to build a \emph{circular} data structure, and how to let something refer to \emph{itself} in its own definition. A plain \code{let} answers neither, for the reason of \Cref{sec:recproblem}: the right-hand side is evaluated before the name is bound. The recitation shows two ways around this. The first ties the knot with boxes, as in \Cref{sec:cycles}. The second uses \emph{fixed-point combinators}, which give recursion with no mutation at all. A challenge at the end uses one of them to implement \code{rec} in the interpreter. The code is plain \code{\#lang racket}, with tests written with \code{check-expect} from \code{test-engine/racket-tests}, except for the challenge, which is \code{plai-typed}.

\subsection{Tying the Knot with Boxes}

\code{self-box} is a box that contains itself, and \code{two-cycle} is two boxes that point to each other, $\code{b1} \to \code{b2} \to \code{b1}$:
\begin{lstlisting}
(define self-box
  (let ([b (box 'dummy)])
    (begin
      (set-box! b b)
      b)))
(define two-cycle
  (let* ([b1 (box 'dummy)]
         [b2 (box b1)])
    (begin
      (set-box! b1 b2)
      b1)))

(check-expect (eq? (unbox self-box) self-box) #t)
(check-expect (eq? (unbox (unbox two-cycle)) two-cycle) #t)
(check-expect (eq? (unbox two-cycle) two-cycle) #f)
\end{lstlisting}
The tests use \code{eq?}, which asks whether two values are the \emph{same} object in memory. Unboxing \code{self-box} once gets back to it; \code{two-cycle} needs two steps, and one step gets to the other box. Racket prints the two values as \code{\#0='\#\&\#0\#} and \code{\#0='\#\&\#\&\#0\#} (\Cref{sec:cycles}).

\begin{exercise}[Recursive functions with boxes]
Implement factorial as \code{fact-box} by putting the function in a box, without letting \code{fact-box} call itself by name.
\end{exercise}
\begin{solution}
This is the factorial of \Cref{sec:cycles}, with the call \code{((unbox fact) n)} at the end:
\begin{lstlisting}
(define (fact-box n)
  (let* ([fact (box 'dummy)]
         [fact-fun (lambda (k)
                     (if (equal? k 0)
                         1
                         (* k ((unbox fact) (- k 1)))))])
    (begin
      (set-box! fact fact-fun)
      ((unbox fact) n))))

(check-expect (fact-box 0) 1)
(check-expect (fact-box 5) 120)
\end{lstlisting}
\end{solution}

\begin{pitfall}[\code{(unbox fact)} is the function, not the call]
\code{(* k (unbox fact) (- k 1))} is also syntactically fine, but it multiplies three things: \code{k}, the function itself, and \code{k - 1}. It fails with \code{*: contract violation, expected: number?, given: \#<procedure:fact-fun>}. The test \code{(fact-box 0)} still passes, because the base case never reaches the multiplication, so only \code{(fact-box 5)} shows the bug. The call needs its own parentheses: \code{((unbox fact) (- k 1))}.
\end{pitfall}

\subsubsection{A loop with no recursive definition} What does this program do?
\begin{lstlisting}
(let ([nonterminating (box 'dummy)])
  (begin
    (set-box! nonterminating (lambda (x) ((unbox nonterminating) x)))
    ((unbox nonterminating) 5)))
\end{lstlisting}
It runs forever, although nothing in it is defined in terms of itself. The loop is in the store, not in the program text. The \code{lambda} reaches the box through its environment, and after \code{set-box!} the box holds that same \code{lambda}. So each call unboxes the function and calls it again. \code{(define (loop x) (loop x))} loops for the same reason: a recursive \code{define} also ties a knot, and Racket ties it for us (\Cref{sec:recboxes}). The call is in tail position, so Racket does not grow the stack; memory use stays flat and the program simply never stops.

\begin{exercise}[Circular linked lists]
Write \code{circ-list}, which takes $n \ge 1$ and returns a circular list of $n$ boxes. Each box holds a pair \code{(cons value next)}. The head holds $n$, the next box $n - 1$, and so on down to 1, and the box holding 1 points back to the head:
\[
\text{head } (n) \to (n-1) \to \dots \to (1) \to \text{head}
\]
For $n < 1$, signal \code{(error 'circ-list "n must be at least 1")}. Use a recursive helper stored in a box, which takes a number \code{k} and a box \code{b}, fills \code{b} with \code{k} and the next box, and builds the rest of the list.
\end{exercise}
\begin{solution}
Two knots are tied here. The head is allocated first as a placeholder, so the last box can point to it. And the helper \code{build} is recursive, so it lives in a box too:
\begin{lstlisting}
(define (circ-list n)
  (if (< n 1)
      (error 'circ-list "n must be at least 1")
      (let* ([head (box 'dummy)]
             [build (box 'dummy)]
             [build-fun
              (lambda (k b)
                (if (equal? k 1)
                    (set-box! b (cons 1 head))
                    (let ([next-box (box 'dummy)])
                      (begin
                        (set-box! b (cons k next-box))
                        ((unbox build) (- k 1) next-box)))))])
        (begin
          (set-box! build build-fun)
          ((unbox build) n head)
          head))))
\end{lstlisting}
\code{build} fills the box it is given, which is the head on the first call, and allocates the next box only when $k > 1$. The box for 1 is filled with \code{head} instead of a new box, which closes the cycle. With the helpers
\begin{lstlisting}
(define (node-value b) (car (unbox b)))
(define (node-next b) (cdr (unbox b)))
\end{lstlisting}
the expected results are
\begin{lstlisting}
(define circ-1 (circ-list 1))
(define circ-3 (circ-list 3))
(check-expect (node-value circ-3) 3)
(check-expect (node-value (node-next circ-3)) 2)
(check-expect (eq? (node-next circ-1) circ-1) #t)
(check-expect (eq? (node-next circ-3) circ-3) #f)
(check-expect (eq? (node-next (node-next (node-next circ-3))) circ-3) #t)
(check-error (circ-list 0) "circ-list: n must be at least 1")
\end{lstlisting}
\code{circ-1} is a one-box cycle, like \code{self-box}, so its next box is itself. Racket prints \code{(circ-list 3)} as \code{\#0='\#\&(3 . \#\&(2 . \#\&(1 . \#0\#)))}.
\end{solution}

\begin{exercise}[Recursive trees]
A tree is \code{(leaf $v$)} or \code{(node $t_1$ $t_2$)}, written as an s-expression \code{'(leaf v)} or \code{'(node left right)}. Write \code{tree-box}, which turns a nested list into a tree: a number becomes a leaf, \code{(list l r)} becomes a node, and anything else is an error \code{"tree: not a tree"}. Tie the knot with a box; \code{tree-box} must not call itself by name.
\begin{center}
\begin{tabular}{@{}ll@{}}
\toprule
input & output \\
\midrule
\code{5} & \code{'(leaf 5)} \\
\code{(list 5 (list 6 7))} & \code{'(node (leaf 5) (node (leaf 6) (leaf 7)))} \\
\code{(list 1 2 3)} & error \\
\bottomrule
\end{tabular}
\end{center}
\end{exercise}
\begin{solution}
The pattern \code{(? number? n)} matches any number and binds it to \code{n}. A node needs two recursive calls, one for each subtree:
\begin{lstlisting}
(define (tree-box input)
  (let* ([tree (box 'dummy)]
         [tree-fun
          (lambda (t)
            (match t
              [(? number? n) (list 'leaf n)]
              [(list left right) (list 'node ((unbox tree) left) ((unbox tree) right))]
              [_ (error 'tree "not a tree")]))])
    (begin
      (set-box! tree tree-fun)
      ((unbox tree) input))))

(check-expect (tree-box 5) '(leaf 5))
(check-expect (tree-box (list 5 (list 6 7)))
              '(node (leaf 5) (node (leaf 6) (leaf 7))))
(check-error (tree-box (list 1 2 3)) "tree: not a tree")
(check-error (tree-box (list 5 (list 6))) "tree: not a tree")
\end{lstlisting}
\code{(list 5 (list 6))} is an error because the inner list has one element, which matches neither a number nor a two-element list. The pitfall above applies here as well: \code{(list 'node ((unbox tree) left) (unbox tree) right)} builds a list of four elements, with the function itself as the third one.
\end{solution}

\subsection{Recursion without Mutation}

Can we get recursion without boxes? The idea is that a function does not refer to itself. Instead, it takes ``the recursive call'' as an argument, as the function passed to \code{fix} in \Cref{sec:cycles} does, but one argument at a time:
\begin{lstlisting}
(define fact-gen
  (lambda (rec)
    (lambda (n)
      (if (equal? 0 n)
          1
          (* n (rec (- n 1)))))))
\end{lstlisting}
\code{fact-gen} is not factorial. Given a function \code{rec} that computes factorial, it returns a function that computes factorial too, by doing one step and leaving the rest to \code{rec}. What we need is a function that is not changed by \code{fact-gen}.

\begin{definition}[Fixed point]
A \emph{fixed point} of a function $f$ is a value $v$ with $f(v) = v$. A \emph{fixed-point combinator} is a function $Y$ such that, for every $f$,
\[
Y\ f = f\ (Y\ f),
\]
that is, $Y\ f$ is a fixed point of $f$.
\end{definition}
If such a $Y$ exists, \code{(Y fact-gen)} is factorial, because unfolding the equation once gives exactly the factorial step:
\begin{lstlisting}[style=repl]
Y fact-gen
=  fact-gen (Y fact-gen)
=  (lambda (n) (if (equal? 0 n) 1 (* n ((Y fact-gen) (- n 1)))))
\end{lstlisting}
Here \code{=} means the two sides are equal, but not by a single evaluation step; \code{=>} below is one step.

\subsubsection{The Y combinator} One such $Y$ is
\begin{lstlisting}
(define Y
  (lambda (h)
    ((lambda (x) (h (x x)))
     (lambda (x) (h (x x))))))
\end{lstlisting}
The trick is \emph{self-application}, \code{(x x)}. The same function is written twice, and the first copy is applied to the second. So inside its body, \code{x} is that function itself, and \code{(x x)} rebuilds the whole application. \code{Y} satisfies the equation:
\begin{lstlisting}[style=repl]
Y f
=> (lambda (x) (f (x x))) (lambda (x) (f (x x)))
=> f ((lambda (x) (f (x x))) (lambda (x) (f (x x))))
=  f (Y f)
\end{lstlisting}
The last step is \code{=}, not \code{=>}: the argument of \code{f} is what \code{(Y f)} stepped to on the first line, but nothing steps back to \code{(Y f)}.

\subsubsection{Why \code{Y} loops in Racket} \code{(define fact-y (Y fact-gen))} never finishes. Write $W$ for \code{(lambda (x) (fact-gen (x x)))}. Then
\begin{lstlisting}[style=repl]
(Y fact-gen)
=> (W W)
=> (fact-gen (W W))
=> (fact-gen (fact-gen (W W)))
=> ...
\end{lstlisting}
Racket evaluates arguments eagerly, so before it can call \code{fact-gen}, it must evaluate the argument \code{(W W)}. That is the expression it started from, one step later. \code{fact-gen} is never called, although it would not even use \code{rec} right away: it would return a \code{lambda} at once.

In a lazy language, \code{Y} works. The argument \code{(W W)} is passed unevaluated, as in the lazy substitution of \Cref{sec:lazy}. \code{fact-gen} returns its \code{(lambda (n) ...)}, and \code{(W W)} is evaluated only when the body calls \code{rec}, which unfolds one more level. A base case that does not call \code{rec} stops the unfolding.

\begin{exercise}[The Z combinator]
Write \code{Z}, the version of \code{Y} for eager evaluation. For every \code{f} and value \code{v} it satisfies
\[
Z\ f\ v = f\ (Z\ f)\ v.
\]
Hint: start from \code{Y}, and delay the self-application by wrapping it in a \code{lambda}: \code{(x x)} becomes \code{(lambda (v) ((x x) v))}.
\end{exercise}
\begin{solution}
This is the $\eta$-expansion of \Cref{sec:recboxes}, applied to both copies of \code{(x x)}:
\begin{lstlisting}
(define Z
  (lambda (h)
    ((lambda (x) (h (lambda (v) ((x x) v))))
     (lambda (x) (h (lambda (v) ((x x) v)))))))

(define fact-z
  (lambda (n) ((Z fact-gen) n)))

(check-expect (fact-z 0) 1)
(check-expect (fact-z 5) 120)
\end{lstlisting}
Write $W'$ for \code{(lambda (x) (fact-gen (lambda (v) ((x x) v))))}. Now
\begin{lstlisting}[style=repl]
(Z fact-gen)
=> (W' W')
=> (fact-gen (lambda (v) ((W' W') v)))
=> (lambda (n) (if (equal? 0 n) 1 (* n (rec (- n 1)))))
       where rec = (lambda (v) ((W' W') v))
\end{lstlisting}
The argument of \code{fact-gen} is a \code{lambda}, which is already a value, so \code{fact-gen} is called, and the result is a function. Evaluation stops there. \code{(W' W')} is evaluated only when \code{rec} is called, and then it produces the same function again, one level down:
\begin{lstlisting}[style=repl]
((Z fact-gen) 3)
=> (* 3 (rec 2))
=> (* 3 ((W' W') 2))
=> (* 3 (* 2 (rec 1)))
=> (* 3 (* 2 (* 1 (rec 0))))
=> (* 3 (* 2 (* 1 1)))
=> 6
\end{lstlisting}
Each line skips the steps that unfold \code{(W' W')} into the factorial function again.
\end{solution}

\begin{exercise}[Putting it together]
Rewrite \code{tree-box} with \code{Z} instead of a box. First write \code{tree-gen}, which takes the recursive call as an argument, like \code{fact-gen}.
\end{exercise}
\begin{solution}
The body of \code{tree-fun} is unchanged, except that the recursive call is the argument \code{rec} instead of \code{(unbox tree)}:
\begin{lstlisting}
(define tree-gen
  (lambda (rec)
    (lambda (t)
      (match t
        [(? number? n) (list 'leaf n)]
        [(list left right) (list 'node (rec left) (rec right))]
        [_ (error 'tree "not a tree")]))))
(define tree-z
  (lambda (t) ((Z tree-gen) t)))
\end{lstlisting}
The four tests of \code{tree-box} pass for \code{tree-z} as well.
\end{solution}

\subsubsection{\code{Z} versus \code{fix}} Both turn a function like \code{fact-gen} into a recursive function. \code{fix} ties the knot with a box: there is one \code{frec}, and it finds itself in the store. \code{Z} uses nothing but functions: the function never finds ``itself'', it rebuilds a copy with \code{(x x)} at every recursive call. (The Racket \code{fix} of \Cref{sec:cycles} takes \code{rec} and \code{n} together, \code{(lambda (rec n) ...)}; \code{Z} takes them one at a time, like the \code{fix} of \Cref{sec:recboxes}.)

\subsection{Challenge: \texttt{rec} via \texttt{Z}}
\label{sec:rec7-interp}

The file \code{recitation7-interp.rkt} contains the interpreter of \Cref{sec:recproblem} with \code{if0C} and \code{recC}, but with no \code{recV} and no boxes: \tyname{Value} has only \code{numV} and \code{closV}. The task is to implement \code{rec} by rewriting it into a call to \code{Z}, as \Cref{sec:recboxes} rewrote it into a call to \code{fix}. A helper \code{run} parses and evaluates:
\begin{lstlisting}
(define (run (s : s-expression)) : Value
  (eval (parse s)))
\end{lstlisting}

\begin{exercise}[Z in our language]
Write \code{Z} as a program in our language. \code{lambda} takes exactly one argument, written \code{(lambda x body)}, and application is \code{(e1 e2)}.
\end{exercise}
\begin{solution}
The Racket \code{Z} already uses one-argument functions only, so only the syntax changes:
\begin{lstlisting}
(define Z
  (parse '(lambda h
            ((lambda x (h (lambda v ((x x) v))))
             (lambda x (h (lambda v ((x x) v))))))))

(define fact-gen
  (parse '(lambda self
            (lambda n
              (if0 n 1 (* n (self (+ n -1))))))))

(test (eval (appC (appC Z fact-gen) (numC 0))) (numV 1))
(test (eval (appC (appC Z fact-gen) (numC 5))) (numV 120))
\end{lstlisting}
In \code{fact-gen}, which the handout provides, the parameter is called \code{self} and not \code{rec}, because \code{rec} is a keyword of our language: \code{(rec (+ n -1))} would parse as a \code{rec} expression.
\end{solution}

\begin{exercise}[\code{rec}]
Fill in the \code{recC} case of \code{eval-env}.
\end{exercise}
\begin{solution}
Exactly as with \code{fix}, \code{(rec f x e)} becomes \code{(Z (lambda f (lambda x e)))}:
\begin{lstlisting}
    [recC (f x e) (eval-env env (appC Z (lambdaC f (lambdaC x e))))]
\end{lstlisting}
\code{(lambda f (lambda x e))} plays the role of \code{fact-gen}. All four tests pass:
\begin{lstlisting}
(test (run '((rec fact x (if0 x 1 (* x (fact (+ x -1))))) 5)) (numV 120))
(test (run '(let fact (rec f x (if0 x 1 (* x (f (+ x -1)))))
              (fact 6)))
      (numV 720))
(test (run '(let fibb (rec f x (if0 x
                                    0
                                    (if0 (+ x -1)
                                         1
                                         (+ (f (+ x -1)) (f (+ x -2))))))
              (fibb 6)))
      (numV 8))
(test (run '(let k 10
              (let f (rec g x (if0 x k (g (+ x -1))))
                (f 3))))
      (numV 10))
\end{lstlisting}
The last test checks that the body of a \code{rec} can still use variables from outside it. It passes for the reason of \Cref{sec:recboxes}: \code{Z} is a closed program, and the user's \code{e} is passed to it as an argument, so \code{Z}'s names \code{h}, \code{x}, and \code{v} never capture anything in \code{e}. That holds even though the tests use \code{x} as their own parameter name.
\end{solution}

\begin{exercise}[Y instead of Z]
Replace \code{Z} by \code{Y} in the \code{recC} case and run the file again. What happens, and why?
\end{exercise}
\begin{solution}
The \code{Z} tests still pass, but the first \code{rec} test never finishes, and memory use grows steadily until DrRacket is stopped. This is the loop of the Racket \code{Y}, one level down. Our language is eager too: \code{appC} evaluates the argument \code{e2} to a value before it runs the body. So \code{(h (x x))} evaluates \code{(x x)}, which evaluates \code{(h (x x))} again, without ever calling \code{h}. Unlike the Racket loop above, this one also grows the Racket stack: each \code{appC} evaluates its argument with a nested call to \code{eval-env}, which is not a tail call.
\end{solution}

\begin{exercise}[Three ways to implement \code{rec}]
For each of \code{recV}, \code{fix}, and \code{Z}: does it need a new kind of \tyname{Value}? Does it need boxes (a store)?
\end{exercise}
\begin{solution}
\begin{center}
\begin{tabular}{@{}>{\raggedright}p{0.2\textwidth}>{\raggedright}p{0.24\textwidth}>{\raggedright}p{0.22\textwidth}>{\raggedright\arraybackslash}p{0.24\textwidth}@{}}
\toprule
& \code{recV} (\Cref{sec:recproblem}) & \code{fix} (\Cref{sec:recboxes}) & \code{Z} (this section) \\
\midrule
new kind of \tyname{Value} & yes, \code{recV} & no & no \\
boxes / a store & no & yes & no \\
interpreter changes & \code{appC} gets a case & only \code{recC} & only \code{recC} \\
where the self-reference is & environment, rebound by each call & a cycle through the store & rebuilt by \code{(x x)} at each call \\
\bottomrule
\end{tabular}
\end{center}
\code{recV} changes the interpreter the most, but it is the most direct: the function's own name is bound to the function, at each call. \code{fix} and \code{Z} change nothing but \code{recC}, because both turn \code{rec} into a program the interpreter can already run. \code{fix} needs the language to have boxes. \code{Z} needs only \code{lambda} and application, which is why recursion needs no special support in an untyped language with first-class functions.
\end{solution}

\begin{remark}
\code{Z} relies on \code{(x x)}, and this self-application cannot be typed directly in a language like \code{plai-typed}: \code{x} would need a type of functions whose argument has that same type. This is the same obstacle as the box that contains itself in \Cref{sec:cycles}. (A recursive \code{define-type} that wraps \code{x} gets around it.) Typed languages therefore usually build recursion in, as \code{define} and \code{letrec} are built into Racket.
\end{remark}

